\documentclass[11pt,a4paper]{article}

% --- Packages ---
\usepackage[utf8]{inputenc}
\usepackage[T1]{fontenc}
\usepackage{lmodern}
\usepackage[margin=1in]{geometry}
\usepackage{amsmath,amssymb,amsthm}
\usepackage{booktabs}
\usepackage{multirow}
\usepackage{graphicx}
\usepackage{xcolor}
\usepackage{listings}
\usepackage{hyperref}
\usepackage{cleveref}
\usepackage{enumitem}
\usepackage{caption}
\usepackage{subcaption}
\usepackage{tabularx}
\usepackage{array}

% --- Listings config ---
\definecolor{leanblue}{RGB}{0,51,153}
\definecolor{leangreen}{RGB}{0,128,0}
\definecolor{leangray}{RGB}{128,128,128}
\definecolor{bglight}{RGB}{248,248,248}

\lstdefinelanguage{Lean4}{
  keywords={def,theorem,lemma,structure,inductive,where,if,then,else,match,with,let,have,by,exact,unfold,split,cases,induction,rename_i,simp,omega,dsimp,rw,fun,some,none,true,false,Prop,Bool,Nat,String,List,Option,deriving},
  sensitive=true,
  morecomment=[l]{--},
  morecomment=[n]{/-}{-/},
  morestring=[b]",
  keywordstyle=\color{leanblue}\bfseries,
  commentstyle=\color{leangreen}\itshape,
  stringstyle=\color{red!60!black},
}

\lstset{
  basicstyle=\ttfamily\small,
  backgroundcolor=\color{bglight},
  frame=single,
  framerule=0.4pt,
  numbers=left,
  numberstyle=\tiny\color{leangray},
  breaklines=true,
  captionpos=b,
  tabsize=2,
  xleftmargin=2em,
  framexleftmargin=1.5em,
}

% --- Hyperref config ---
\hypersetup{
  colorlinks=true,
  linkcolor=blue!60!black,
  citecolor=blue!60!black,
  urlcolor=blue!60!black,
}

% --- Title ---
\title{Formally Verified Harness Engineering:\\
  Using Lean~4 to Prove Correctness Properties\\
  of LLM Orchestration Code}

\author{Samy Djouder}

\date{April 2026}

\begin{document}
\maketitle

% ============================================================================
\begin{abstract}
% ============================================================================

LLM orchestration frameworks coordinate multi-agent systems through
shared state machines that manage token accounting, budget enforcement,
and execution control flow. Despite widespread adoption, none of these
frameworks provides formal specifications or machine-checked proofs of
its core invariants. We present a three-phase empirical study applied
to an 800-line Python orchestration harness backed by a Lean~4
specification with 15 machine-checked theorems. Phase~1 generates
420 mutants via \texttt{mutmut} to establish a population baseline
(55\% raw, 80\% adjusted kill rate). Phase~2 tests 30 hand-crafted
mutations against 4 verification methods---unit, integration,
property-based, and formal---to measure their complementarity.
Phase~3 derives 11 mutations by negating each conjunct of the Lean~4
well-formedness predicate. We find that formal verification detects
80\% of mutations within its scope but covers only 33\% of the mutation
surface; property-based testing catches 3 mutations no other method
detects; and one spec-guided mutation (replacing cost accumulation
\texttt{+=} with assignment \texttt{=}) passes both the Lean~4 spec
and all 65 tests, exposing a missing cost-monotonicity invariant.
Structural analysis of smolagents and pydantic-ai confirms that 8 of
11 gap categories are present in all three frameworks. All proofs
compile without \texttt{sorry}; all artifacts are publicly available.

\end{abstract}

% ============================================================================
\section{Introduction}
\label{sec:intro}
% ============================================================================

The adoption of large language models (LLMs) has driven the development
of orchestration frameworks---software systems that coordinate multiple
LLM-powered agents through shared state, tool registries, and control
flow graphs~\cite{langgraph2024,crewai2024,autogen2024}. These
frameworks implement invariants that matter in production: token counts
must be tracked accurately for billing, budget limits must halt
execution before overspend, and termination must be guaranteed to
prevent runaway agent loops.

No production LLM orchestration framework provides formal
specifications or machine-checked proofs of these properties. The
existing test suites for LangGraph, smolagents, and pydantic-ai
exercise functional behavior but do not verify state machine
invariants. A subtle off-by-one in a budget check or a missing cost
accumulation can lead to unbounded API spending---a failure mode that
unit tests are unlikely to exercise because it requires specific
combinations of cost values at the budget boundary.

This paper bridges formal verification (Lean~4) and mutation testing
to characterize the \emph{formal/informal boundary}: the set of bugs
that formal methods catch but testing misses, and the (larger) set
that testing catches but formal methods cannot reach.

\paragraph{Contributions.}
\begin{enumerate}[leftmargin=*,nosep]
  \item A Lean~4 specification of 5 well-formedness invariants for
        LLM orchestration state machines, with 15 machine-checked
        theorems covering safety, liveness, and synchronization
        (\Cref{sec:spec}).
  \item A three-phase mutation study comparing 4 verification
        methods across 461 mutations (420 automated + 30 targeted
        + 11 spec-guided, with partial overlap)
        (\Cref{sec:method,sec:results}).
  \item A spec-guided mutation technique that negates formal
        invariants to produce semantically meaningful mutations,
        exposing a concrete abstraction gap in the Lean~4 spec
        (\Cref{sec:specguided}).
  \item A cross-system structural analysis showing that the gap
        taxonomy holds across three independent frameworks
        (\Cref{sec:crosssystem}).
\end{enumerate}

\paragraph{Research questions.}
\begin{description}[leftmargin=*,nosep]
  \item[RQ1] How complementary are formal verification and
    testing for detecting mutations in LLM orchestration code?
  \item[RQ2] What fraction of the mutation surface falls within the
    scope of a Lean~4 specification?
  \item[RQ3] Does the gap taxonomy generalize beyond a single
    implementation?
\end{description}

% ============================================================================
\section{Background}
\label{sec:background}
% ============================================================================

\subsection{LLM Orchestration Frameworks}

An LLM orchestration framework manages the execution loop of one or
more LLM-powered agents. The common architecture includes a
\emph{state machine} tracking step counts, token usage, and costs;
a \emph{tool registry} mapping function names to executable code;
a \emph{routing layer} selecting which agent runs next; and a
\emph{budget/context manager} enforcing resource limits. LangGraph~\cite{langgraph2024},
CrewAI~\cite{crewai2024}, AutoGen~\cite{autogen2024},
smolagents~\cite{smolagents2024}, and
pydantic-ai~\cite{pydanticai2024} all follow this pattern.

\subsection{Mutation Testing}

Mutation testing evaluates test effectiveness by injecting small
syntactic faults (mutants) into source code and checking whether tests
detect them~\cite{demillo1978hints,jia2011analysis}. A mutant
\emph{killed} by a test confirms the test exercises the mutated
behavior; a \emph{surviving} mutant may indicate a test gap.
The mutation score (killed/total) is the standard adequacy metric,
though equivalent mutants---changes with no observable
effect---inflate the denominator and require
classification~\cite{papadakis2019mutation,ammann2014establishing}.
Petrovi\'c and Ivankovi\'c~\cite{petrovic2018state} describe
deployment of mutation testing at Google scale.

\subsection{Lean~4}

Lean~4 is a theorem prover and programming language with dependent
types, tactic-based proofs, and a compiled
runtime~\cite{demoura2021lean4}. We use it to state specifications
as propositions and verify them with constructive proofs that the
kernel checks exhaustively. A proven theorem holds for all inputs
within the model's scope---unlike testing, which covers only the
exercised inputs.

% ============================================================================
\section{System Under Study}
\label{sec:system}
% ============================================================================

Our subject system, \textsc{excelsior-harness}, is an 8-module
Python harness (780~LOC) built to the standard architecture
described in \Cref{sec:background}:

\begin{itemize}[nosep]
  \item \texttt{state.py} (73~LOC): Immutable \texttt{AgentState}
    with functional updates (\texttt{add\_message},
    \texttt{record\_usage}, \texttt{checkpoint}, \texttt{terminate}).
  \item \texttt{budget.py} (76~LOC): Per-model cost tracking with
    hard USD ceiling and \texttt{BudgetExceeded} exception.
  \item \texttt{tools.py} (138~LOC): Decorator-based tool
    registration with JSON schema extraction from type hints.
  \item \texttt{agents.py} (166~LOC): Role-based agents
    (supervisor, worker) with per-agent tool permissions.
  \item \texttt{graph.py} (112~LOC): Directed graph with
    conditional edges for routing.
  \item \texttt{context.py} (97~LOC): Tiktoken-based token counting
    with truncation (keep system message + recent messages).
  \item \texttt{orchestrator.py} (201~LOC): Main loop with
    budget/step checks and exponential-backoff retry.
  \item \texttt{\_types.py} (47~LOC): Shared enums and type aliases.
\end{itemize}

The harness uses mock LLM clients (no real API calls) for
reproducibility. The test suite has 65 pytest tests: 40 unit,
5 integration, and 20 property-based
(Hypothesis~\cite{hypothesis2024}).

% ============================================================================
\section{Lean~4 Specification}
\label{sec:spec}
% ============================================================================

We model the orchestrator as a state machine with explicit step,
budget, token, and checkpoint counters.
\Cref{lst:lean-core} shows the state structure and transition
function (simplified for presentation; the full specification uses
anonymous constructors and is available in the repository).

\begin{lstlisting}[language=Lean4,caption={Orchestrator state and transition function (simplified from repository source).},label=lst:lean-core]
structure OState where
  step : Nat; maxSteps : Nat
  cost : Nat; maxBudget : Nat
  pTok : Nat; cTok : Nat; tTok : Nat
  agent : String; done : Bool
  reason : Option TerminationReason
  ckpts : Nat

def execStep (s : OState) (c : CallResult) (r : Route) : OState :=
  if s.step >= s.maxSteps then            -- PATH 1: max_steps
    { s with done := true, reason := some .max_steps }
  else if s.cost + c.cost > s.maxBudget then -- PATH 2: budget
    { s with step := s.step + 1, cost := s.cost + c.cost,
      pTok := s.pTok + c.pTok, cTok := s.cTok + c.cTok,
      tTok := s.tTok + c.pTok + c.cTok,
      done := true, reason := some .budget_exceeded }
  else match r with
    | .finished => { ... }                -- PATH 3: task_complete
    | .next ag => { ... }                 -- PATH 4: continue
\end{lstlisting}

\subsection{Well-Formedness Predicate}

The specification is organized around a well-formedness predicate
\texttt{wf} that conjoins five invariants:

\begin{lstlisting}[language=Lean4,caption={Well-formedness predicate.},label=lst:wf]
def wf (s : OState) : Prop :=
  s.tTok = s.pTok + s.cTok                                 -- P1
  /\ (s.reason != some .budget_exceeded -> s.cost <= s.maxBudget) -- P2
  /\ s.ckpts <= s.step                                      -- P3
  /\ s.step <= s.maxSteps + 1                                -- P4
  /\ (s.reason = some .budget_exceeded -> s.done = true)     -- P5
\end{lstlisting}

\noindent
\textbf{P1} (Token accounting): Total tokens equal prompt plus
completion. \textbf{P2} (Budget safety): Cost does not exceed budget
unless the system has already flagged budget exhaustion. \textbf{P3}
(Checkpoint monotonicity): Checkpoint count never exceeds step count.
\textbf{P4} (Step bound): Step count is bounded by
$\texttt{maxSteps} + 1$. \textbf{P5} (Termination consistency):
Budget exhaustion implies the done flag is set.

P5 prevents a subtle interaction between the max-steps and
budget-exceeded paths: without it, a state could carry
\texttt{reason = budget\_exceeded} while \texttt{done = false},
breaking downstream logic that checks \texttt{done} to decide
whether to continue the loop.

\subsection{Proven Theorems}

\Cref{tab:theorems} lists the 14 theorems plus a capstone
composition. All are machine-checked with no \texttt{sorry} and no
axioms beyond the standard library.

\begin{table}[htbp]
\centering
\caption{Machine-checked theorems.}
\label{tab:theorems}
\small
\begin{tabular}{@{}llp{6cm}@{}}
\toprule
\textbf{\#} & \textbf{Theorem} & \textbf{Property} \\
\midrule
1 & \texttt{step\_preserves\_wf} & All 5 invariants survive one step ($4 \times 5 = 20$ proof obligations) \\
2 & \texttt{loop\_budget\_safe} & Budget ceiling across full loop (induction on fuel) \\
3 & \texttt{loop\_terminates} & Loop terminates within \texttt{maxSteps} (bounded liveness) \\
4 & \texttt{loop\_tokens} & Token accounting preserved across loop \\
5 & \texttt{tool\_isolation} & Per-agent tool permissions enforced \\
6 & \texttt{loop\_ckpts} & Checkpoints $\leq$ steps across loop \\
\midrule
7 & \texttt{step\_done\_at\_max} & Max-steps $\Rightarrow$ done \\
8 & \texttt{step\_inc\_or\_done} & Step increments or execution terminates \\
9 & \texttt{step\_maxSteps\_eq} & \texttt{maxSteps} is constant through execution \\
10 & \texttt{step\_preserves\_tokens} & Token identity preserved per step \\
11 & \texttt{step\_ckpt\_bounded} & Checkpoint bound per step \\
12 & \texttt{continue\_adds\_ckpt} & Routing to next agent increments checkpoint \\
13 & \texttt{max\_steps\_no\_ckpt} & Max-steps termination does not checkpoint \\
14 & \texttt{no\_perms\_no\_tools} & No permissions entry $\Rightarrow$ no tool access \\
\midrule
C & \texttt{all\_invariants\_preserved} & Capstone: tokens $\wedge$ budget $\wedge$ checkpoints \\
\bottomrule
\end{tabular}
\end{table}

% ============================================================================
\section{Study Design}
\label{sec:method}
% ============================================================================

The study has three phases (\Cref{fig:phases}). Phase~1 gives
unbiased population statistics but can only test against the pytest
suite. Phase~2 enables multi-method comparison but is vulnerable to
selection bias. Phase~3 derives mutations from the formal spec itself
to probe the formal/informal boundary. The phases are complementary
by design: Phase~1 validates that Phase~2's selection is
representative, and Phase~3 targets exactly the behaviors that the
Lean spec claims to verify.

\begin{figure}[htbp]
\centering
\begin{tabular}{@{}ccc@{}}
\fbox{\parbox{3.8cm}{\centering\textbf{Phase 1}\\420 mutmut mutants\\1 test suite\\Population baseline}} &
$\longrightarrow$ &
\fbox{\parbox{3.8cm}{\centering\textbf{Phase 2}\\30 hand-crafted\\4 methods\\Multi-method comparison}} \\[1.5em]
& & $\downarrow$ \\[0.5em]
& &
\fbox{\parbox{3.8cm}{\centering\textbf{Phase 3}\\11 spec-guided\\Lean negation\\Boundary analysis}} \\
\end{tabular}
\caption{Three-phase study design. Phase~1 establishes baselines;
Phase~2 compares methods; Phase~3 probes the specification boundary.}
\label{fig:phases}
\end{figure}

\subsection{Phase 1: Automated Baseline}

\texttt{mutmut}~2.5.1~\cite{mutmut2024} generated 420 AST-level
mutants across all 8 source files. Each mutant was tested against the
65-test suite. Survivors were classified into four tiers following
Papadakis et al.~\cite{papadakis2019mutation}:

\begin{enumerate}[nosep]
  \item \textbf{Equivalent} (76): No observable behavior change
    (logging strings, enum repr methods, error message text).
  \item \textbf{Test infrastructure} (12): Changes to
    \texttt{MockResponse} or \texttt{MockLLMClient} defaults.
  \item \textbf{Configuration} (40): Cost table rates, retry
    parameters, default buffer sizes---tuning constants with no test
    coverage by design.
  \item \textbf{Real gap} (60): Mutations to arithmetic, control
    flow, or state management in production code paths.
\end{enumerate}

We count the 8 ``suspicious'' mutants (timeout during testing)
as killed, following \texttt{mutmut}'s convention. This gives
232 killed, 188 survived.

\subsection{Phase 2: Targeted Multi-Method Comparison}

30 mutations were selected to cover all 4 semantic categories
(logic, boundary, missing\_check, state) and all 8 source files.
13 of 30 have direct Phase~1 counterparts (e.g., M10 $\leftrightarrow$
mutant~\#69, M02 $\leftrightarrow$ mutant~\#64), mitigating selection
bias. Each mutation was tested against 4 methods:

\begin{enumerate}[nosep]
  \item \textbf{Unit tests}: 40 pytest tests across 7 test files.
  \item \textbf{Integration tests}: 5 end-to-end orchestration runs.
  \item \textbf{Property-based tests}: 20 Hypothesis tests
    (13 mirror Lean invariants, 7 cover domain-specific properties).
  \item \textbf{Lean~4 formal spec}: 10 of 30 mutations have spec
    equivalents; detection via \texttt{lake build} pass/fail.
\end{enumerate}

We also planned LLM-based code review as a 5th method but could not
run it (no API key available during the study). We exclude it from
all analysis rather than report zeros.

\subsection{Phase 3: Spec-Guided Mutation Generation}
\label{sec:specguided-method}

For each conjunct $P_i$ of the \texttt{wf} predicate, we identify
the corresponding Lean expression and Python implementation, then
construct 2--3 minimal mutations that negate or weaken $P_i$:

\begin{itemize}[nosep]
  \item \textbf{P1}: Drop completion tokens from total, or
    double-count prompt tokens.
  \item \textbf{P2}: Invert budget comparison, weaken to off-by-one,
    or replace accumulation (\texttt{+=}) with assignment (\texttt{=}).
  \item \textbf{P3}: Double checkpoint increment, or add a spurious
    checkpoint on the max-steps path.
  \item \textbf{P4}: Double or skip the step increment.
  \item \textbf{P5}: Clear the done flag on budget-exceeded, or
    assign the wrong termination reason.
\end{itemize}

This produces 11 mutations, each with a clear semantic interpretation
tied to a specific safety property. Each is tested against Lean
(\texttt{lake build}), the full pytest suite, and the property-based
tests independently.

% ============================================================================
\section{Results}
\label{sec:results}
% ============================================================================

\subsection{Phase 1: Automated Baseline}

\Cref{tab:phase1} summarizes the results.

\begin{table}[htbp]
\centering
\caption{Phase 1 results (mutmut, 420 mutants). Confidence intervals
are Wilson score intervals~\cite{wilson1927probable}.}
\label{tab:phase1}
\small
\begin{tabular}{@{}lrl@{}}
\toprule
\textbf{Metric} & \textbf{Value} & \textbf{95\% Wilson CI} \\
\midrule
Total mutants & 420 & \\
Killed (incl.\ suspicious) & 232 & \\
Survived & 188 & \\
Raw mutation score & 55\% & [50\%, 60\%] \\
Adjusted (excl.\ equiv/config/infra) & 79\% & [74\%, 84\%] \\
\midrule
Real gaps & 60 & \\
Equivalent & 76 & \\
Test infrastructure & 12 & \\
Configuration & 40 & \\
\bottomrule
\end{tabular}
\end{table}

\Cref{tab:byfile} breaks scores down by file. \texttt{state.py}
scores highest (94\%), which is expected: it is the module most
directly modeled by the Lean spec and thus best tested.
\texttt{agents.py} scores lowest (45\%), reflecting the 21 real
gaps in message construction and tool permission logic---behaviors
entirely outside the formal model.

\begin{table}[htbp]
\centering
\caption{Phase 1 scores by source file.}
\label{tab:byfile}
\small
\begin{tabular}{@{}lrrrr@{}}
\toprule
\textbf{File} & \textbf{Total} & \textbf{Killed} & \textbf{Real gaps} & \textbf{Score} \\
\midrule
\texttt{state.py} & 48 & 45 & 2 & 94\% \\
\texttt{graph.py} & 42 & 25 & 1 & 60\% \\
\texttt{context.py} & 58 & 34 & 7 & 59\% \\
\texttt{orchestrator.py} & 66 & 38 & 11 & 58\% \\
\texttt{tools.py} & 55 & 30 & 11 & 55\% \\
\texttt{budget.py} & 51 & 26 & 7 & 51\% \\
\texttt{agents.py} & 76 & 34 & 21 & 45\% \\
\texttt{\_types.py} & 24 & 0 & 0 & 0\% \\
\bottomrule
\end{tabular}
\end{table}

Of 188 survivors, 15 were mapped to the Lean spec. Only 2 fall within
spec scope; Lean catches 1 (mutant~\#69, an off-by-one in the budget
comparison). The other (mutant~\#64, cost $+=$ replaced with $=$) passes
the spec---see \Cref{sec:specguided} for the explanation. The remaining
13 mapped survivors involve rate-level arithmetic, per-model tracking,
loop convergence patterns, and checkpoint data structures that the
Lean model abstracts away.

\subsection{Phase 2: Multi-Method Comparison}

\Cref{tab:phase2} shows detection rates across the 4 methods.

\begin{table}[htbp]
\centering
\caption{Phase 2 detection rates (30 mutations, 4 methods).}
\label{tab:phase2}
\small
\begin{tabular}{@{}lrrrl@{}}
\toprule
\textbf{Method} & \textbf{Detected} & \textbf{Rate} & \textbf{Unique} & \textbf{95\% CI} \\
\midrule
Unit tests & 14/30 & 47\% & 4 & [30\%, 64\%] \\
Integration tests & 6/30 & 20\% & 0 & [10\%, 37\%] \\
Property-based tests & 10/30 & 33\% & 3 & [19\%, 51\%] \\
Lean~4 (all 30) & 8/30 & 27\% & 1 & [14\%, 44\%] \\
Lean~4 (10 in scope) & 8/10 & 80\% & -- & [49\%, 94\%] \\
\midrule
Any method & 18/30 & 60\% & -- & \\
\bottomrule
\end{tabular}
\end{table}

Three results stand out. First, property-based testing provides
3 unique detections---M07 (prompt rate used for both token types),
M13 (cost divided by 100 instead of 1000, causing 10$\times$
inflation), and M26 (returning \texttt{self} instead of a new copy
from \texttt{record\_usage}). These involve numerical invariants
that example-based tests do not exercise with the right input
combinations.

Second, Lean provides 1 unique detection: M10, where the budget
check is weakened from \texttt{>} to \texttt{>=}. None of the
65 tests construct a scenario where cost exactly equals the budget,
so the off-by-one is invisible to testing.

Third, 12 mutations survive all methods. Of these, 11 target
behaviors outside the Lean spec (message construction, context
truncation, tool execution paths, routing logic). The 12th (M06,
swapping prompt and completion token counts) satisfies
$\texttt{tTok} = \texttt{pTok} + \texttt{cTok}$ by commutativity
of addition---a genuine spec weakness, since the identity holds
regardless of which count is which.

\Cref{tab:categories} gives per-category rates.

\begin{table}[htbp]
\centering
\caption{Phase 2 detection by mutation category.}
\label{tab:categories}
\small
\begin{tabular}{@{}lrrrrr@{}}
\toprule
\textbf{Category} & \textbf{n} & \textbf{Unit} & \textbf{Integ.} & \textbf{Prop.} & \textbf{Lean} \\
\midrule
Logic & 8 & 25\% & 12\% & 38\% & 12\% \\
Boundary & 8 & 50\% & 25\% & 25\% & 38\% \\
Missing check & 8 & 62\% & 12\% & 25\% & 25\% \\
State & 6 & 50\% & 33\% & 50\% & 33\% \\
\bottomrule
\end{tabular}
\end{table}

\subsection{Phase 3: Spec-Guided Mutations}
\label{sec:specguided}

\Cref{tab:specguided} presents the 11 spec-guided mutations.

\begin{table}[htbp]
\centering
\caption{Spec-guided mutation results. Each row negates one conjunct of \texttt{wf}.}
\label{tab:specguided}
\small
\begin{tabular}{@{}llp{4.5cm}ccc@{}}
\toprule
\textbf{ID} & \textbf{Inv.} & \textbf{Description} & \textbf{Lean} & \textbf{Tests} & \textbf{Props} \\
\midrule
SGM-P1a & P1 & Drop completion tokens from total & \checkmark & \checkmark & \checkmark \\
SGM-P1b & P1 & Double-count prompt tokens & \checkmark & \checkmark & \checkmark \\
SGM-P2a & P2 & Invert budget check ($>$ to $<$) & \checkmark & \checkmark & \checkmark \\
SGM-P2b & P2 & Weaken budget ($>$ to $\geq$) & \checkmark & $\times$ & $\times$ \\
SGM-P2c & P2 & Replace $\mathrel{+}=$ with $=$ in cost & $\times$ & $\times$ & $\times$ \\
SGM-P3a & P3 & Double checkpoint increment & \checkmark & \checkmark & \checkmark \\
SGM-P3b & P3 & Spurious checkpoint at max-steps & \checkmark & \checkmark & $\times$ \\
SGM-P4a & P4 & Double step increment & \checkmark & \checkmark & \checkmark \\
SGM-P4b & P4 & Skip step increment & \checkmark & \checkmark & \checkmark \\
SGM-P5a & P5 & Budget-exceeded with done=false & \checkmark & \checkmark & $\times$ \\
SGM-P5b & P5 & Wrong termination reason & \checkmark & \checkmark & $\times$ \\
\midrule
\multicolumn{3}{@{}l}{\textbf{Detection rate}} & \textbf{90.9\%} & \textbf{81.8\%} & \textbf{54.5\%} \\
\bottomrule
\end{tabular}
\end{table}

\paragraph{SGM-P2b (Lean-only detection).}
Weakening the budget check from \texttt{>} to \texttt{>=} is caught
by the \texttt{continue\_adds\_ckpt} theorem in Lean, which requires
that cost plus the new call's cost does \emph{not} exceed the budget
for execution to continue. No test constructs a scenario where
accumulated cost exactly equals the budget, so 0 of 65 tests detect
this mutation. This is the same bug as Phase~2 mutation M10 and
Phase~1 mutant~\#69, independently confirming the finding across all
three study phases.

\paragraph{SGM-P2c (the abstraction gap).}
Replacing cost accumulation (\texttt{+=}) with simple assignment
(\texttt{=}) in \texttt{total\_cost} passes both the Lean spec and all
65 tests. The \texttt{wf} predicate checks
$\texttt{cost} \leq \texttt{maxBudget}$ per step but has no
cost-monotonicity invariant. Under natural number arithmetic,
\[
  c.\text{cost} \;\leq\; s.\text{cost} + c.\text{cost}
  \;\leq\; \text{maxBudget},
\]
so dropping the accumulation still satisfies
every theorem. This is a spec weakness: adding a sixth conjunct
$\textit{new}.\texttt{cost} \geq \textit{old}.\texttt{cost}$
to \texttt{wf} would catch it.
The same bug appears independently as Phase~1 mutant~\#64.

\subsection{Summary}

\Cref{tab:combined} collects detection rates across phases.

\begin{table}[htbp]
\centering
\caption{Detection rates across study phases.}
\label{tab:combined}
\small
\begin{tabular}{@{}lrrr@{}}
\toprule
\textbf{Method} & \textbf{Phase 1} & \textbf{Phase 2} & \textbf{Phase 3} \\
 & \textbf{(n=420)} & \textbf{(n=30)} & \textbf{(n=11)} \\
\midrule
Unit tests & \multirow{3}{*}{55\% raw / 79\% adj} & 47\% & -- \\
Integration tests & & 20\% & -- \\
Property-based tests & & 33\% & 54.5\% \\
Lean~4 (in scope) & 50\% (n=2) & 80\% (n=10) & 90.9\% \\
All tests combined & -- & -- & 81.8\% \\
Any method & 55\% & 60\% & -- \\
\bottomrule
\end{tabular}
\end{table}

\paragraph{RQ1.}
Formal verification and testing are complementary. Lean detects
M10/SGM-P2b (off-by-one budget threshold) that no test catches.
Property-based testing detects M07, M13, M26 that Lean cannot model.
Within scope, Lean achieves 80--91\%; outside scope, it provides
nothing. The union of all methods detects 60\% of Phase~2 mutations.

\paragraph{RQ2.}
The Lean specification covers 10 of 30 Phase~2 mutations (33\%)
and 11 of 11 Phase~3 mutations by construction. Of 60 real-gap
survivors in Phase~1, only 2 of 15 mapped mutants fall within spec
scope. The other 87\% involve message construction, tool schemas,
system prompts, context windowing, and routing---none of which the
Lean model addresses.

% ============================================================================
\section{Cross-System Generalizability}
\label{sec:crosssystem}
% ============================================================================

To test whether the gap taxonomy is specific to our harness, we
analyzed the source code of smolagents~1.24.0 (6,384~LOC) and
pydantic-ai~1.77.0 (9,064~LOC). We mapped their modules to the same
7 architectural categories and counted code patterns associated with
each gap category via regex, normalized per 100~LOC. This is a
structural analysis only---we did not run mutation testing on the
external frameworks' test suites.

\begin{table}[htbp]
\centering
\caption{Pattern density (occurrences per 100~LOC) across frameworks.
Measured by regex pattern matching; see \Cref{sec:threats}.}
\label{tab:crosssystem}
\small
\begin{tabular}{@{}lcccc@{}}
\toprule
\textbf{Gap Category} & \textbf{excelsior} & \textbf{smolagents} & \textbf{pydantic-ai} & \textbf{All 3?} \\
\midrule
message\_construction & 3.4 & 1.6 & 4.1 & Yes \\
tool\_schema\_gen & 3.2 & 2.3 & 5.3 & Yes \\
system\_prompt & 1.4 & 1.7 & 2.9 & Yes \\
token\_accounting & 5.5 & 0.6 & 0.6 & Yes \\
context\_window & 2.7 & 0.3 & 0.3 & Yes \\
loop\_termination & 1.5 & 1.1 & 0.4 & Yes \\
checkpoint\_state & 1.5 & 0.3 & 0.1 & Yes \\
retry\_error & 1.5 & 0.1 & 3.1 & Yes \\
\midrule
tool\_permission & 0.5 & 0.0 & 0.1 & Partial \\
cost\_budget & 13.1 & 0.0 & 0.2 & Partial \\
routing\_dispatch & 2.6 & 0.0 & 0.1 & Partial \\
\bottomrule
\end{tabular}
\end{table}

\paragraph{RQ3.}
8 of 11 gap categories appear in all three frameworks
(\Cref{tab:crosssystem}). The 3 partial categories reflect
architectural differences: smolagents has no built-in budget tracking
(it delegates to external monitoring), and pydantic-ai uses implicit
routing rather than an explicit graph. All three frameworks implement
the state machine pattern that our 5 Lean invariants describe, with
identifiable equivalents for each invariant (\Cref{tab:invariantmap}).

\begin{table}[htbp]
\centering
\caption{Invariant equivalents across frameworks.}
\label{tab:invariantmap}
\footnotesize
\begin{tabular}{@{}llll@{}}
\toprule
\textbf{Inv.} & \textbf{excelsior} & \textbf{smolagents} & \textbf{pydantic-ai} \\
\midrule
P1 & state / record\_usage & monitoring / update\_metrics & usage / incr\_usage \\
P2 & budget / BudgetExceeded & monitoring / Monitor & usage / UsageLimits \\
P3 & state / checkpoint & agents / write\_memory & result / stream \\
P4 & orchestrator / max\_steps & agents / max\_steps & agent\_graph / max\_result\_retries \\
P5 & orchestrator / terminate & agents / final\_answer & agent\_graph / check\_for\_result \\
\bottomrule
\end{tabular}
\end{table}

% ============================================================================
\section{Discussion}
\label{sec:discussion}
% ============================================================================

\subsection{The Formal/Informal Boundary}

The three phases converge on a consistent picture of what formal
verification can and cannot do for LLM orchestration code:

\begin{itemize}[nosep]
  \item \textbf{Formal catches what testing misses}: Off-by-one
    boundary conditions (SGM-P2b, M10) where triggering the bug
    requires hitting an exact cost value. No amount of adding unit
    tests will reliably cover this without boundary-value analysis.
  \item \textbf{Testing catches what formal methods cannot reach}:
    Message construction, schema generation, context truncation,
    and routing logic are implementation concerns that the
    mathematical model does not encode.
  \item \textbf{Both miss the same bug}: SGM-P2c exposes a gap where
    the specification and the test suite share a blind spot. The
    Lean model checks cost bounds per-step but not cost monotonicity
    across steps; the tests never accumulate enough cost across
    multiple steps to notice the difference between $+=$ and $=$.
\end{itemize}

The third category is the most useful output of this study: it
identifies a specification weakness that can be fixed by adding one
conjunct to \texttt{wf}, and a testing weakness that can be fixed
by adding a multi-step cost accumulation test.

\subsection{Spec-Guided vs.\ Random Mutations}

The 90.9\% Lean detection rate for spec-guided mutations is higher
than the population baseline, as expected---these mutations target
formally verified invariants by construction. The interesting
mutations are the two that break this pattern:

\begin{itemize}[nosep]
  \item SGM-P2b ($>$ to $\geq$): Lean catches it, but so does
    mutmut mutant~\#69 (Phase~1), confirming it is a real
    population-level gap, not an artifact of our selection.
  \item SGM-P2c ($+=$ to $=$): Neither Lean nor tests catch it.
    Phase~1 mutant~\#64 independently confirms this gap.
\end{itemize}

The cross-phase confirmation matters: both gaps appear in both
the random (Phase~1) and targeted (Phase~3) populations, ruling
out cherry-picking.

\subsection{Practical Implications}

Based on the gap analysis, LLM orchestration frameworks would benefit
from three additions to their test suites:

\begin{enumerate}[nosep]
  \item Property-based tests covering rate arithmetic (M07, M13)
    and copy semantics (M26), where the bug manifests only under
    specific numerical combinations.
  \item Boundary-value tests that set cost exactly equal to the
    budget limit, exercising the $>$ vs.\ $\geq$ edge.
  \item At minimum, runtime assertions for state invariants
    (e.g., asserting cost is non-decreasing after each budget
    recording). Even without a theorem prover, this would have
    caught SGM-P2c.
\end{enumerate}

% ============================================================================
\section{Threats to Validity}
\label{sec:threats}
% ============================================================================

\paragraph{Internal validity.}
Survivor classification in Phase~1 uses rule-based heuristics,
not manual inspection of all 188 survivors. We default to ``real
gap'' for unclassified core-module mutations and report both raw
and adjusted scores. The Lean mapping covers only 15 of 188
survivors; more survivors in \texttt{agents.py}, \texttt{tools.py},
and \texttt{context.py} are unmapped because the spec does not model
those modules.

\paragraph{External validity.}
The harness was built for this study, not extracted from production.
The cross-system analysis (\Cref{sec:crosssystem}) partially
addresses this by showing structural similarity to production
frameworks, but we did not run mutation testing on their test suites.
Pattern counting via regex is a rough proxy for architectural
concerns; false positives are possible.

\paragraph{Construct validity.}
The equivalent mutant problem~\cite{papadakis2019mutation} affects
all mutation studies. Some mutations we classify as equivalent may be
real bugs, and vice versa. 17 of 30 Phase~2 mutations have no
Phase~1 counterpart, which limits our ability to verify that Phase~2
is fully representative.

% ============================================================================
\section{Related Work}
\label{sec:related}
% ============================================================================

\paragraph{Mutation testing.}
Mutation testing was introduced by DeMillo et
al.~\cite{demillo1978hints} and surveyed comprehensively by Jia and
Harman~\cite{jia2011analysis} and Papadakis et
al.~\cite{papadakis2019mutation}. Gopinath et
al.~\cite{gopinath2014mutations} study the difficulty of achieving
high mutation scores; Petrovi\'c and
Ivankovi\'c~\cite{petrovic2018state} report on mutation testing at
Google scale. Our work applies mutation testing to a new domain
(LLM orchestration) and combines it with formal verification to
measure method complementarity.

\paragraph{Formal verification of systems software.}
CompCert~\cite{leroy2009compcert} verified a C compiler;
seL4~\cite{klein2009sel4} verified an OS microkernel. Both
demonstrated that formal verification scales to real systems
with thousands of lines of proof. Our Lean~4 specification is
smaller (319~LOC, 15 theorems) but targets a domain where formal
methods have not been applied before.

\paragraph{Property-based testing.}
\begin{sloppypar}
QuickCheck~\cite{claessen2000quickcheck} pioneered random property-based
testing for Haskell; Hypothesis~\cite{hypothesis2024} brought it to
Python. Our study quantifies their overlap with formal verification:
they share most detections but each finds bugs the other misses---3
unique to property-based tests, 1 unique to Lean.
\end{sloppypar}

\paragraph{LLM agent frameworks.}
LangGraph~\cite{langgraph2024}, CrewAI~\cite{crewai2024},
AutoGen~\cite{autogen2024}, smolagents~\cite{smolagents2024}, and
pydantic-ai~\cite{pydanticai2024} are production frameworks for
multi-agent orchestration. None provides formal specifications.
To our knowledge, this is the first application of theorem proving
to LLM orchestration code.

% ============================================================================
\section{Conclusion}
\label{sec:conclusion}
% ============================================================================

We reported a three-phase mutation study on an LLM orchestration
harness with a Lean~4 specification. The main findings:
formal verification detects 80--91\% of mutations within its
scope, but that scope covers only 13--33\% of the mutation surface;
property-based testing and formal verification are complementary,
each catching bugs the other misses; and spec-guided mutation
generation exposed a concrete specification gap (missing
cost-monotonicity invariant) that neither the formal spec nor the
test suite detected.

The gap taxonomy (8 of 11 categories universal across three
frameworks) suggests these findings are not artifacts of our
specific implementation. The Lean specification and all study
artifacts are available at
\url{https://github.com/[redacted]/excelsior-harness}.

% ============================================================================
% Data Availability
% ============================================================================
\paragraph{Data availability.}
Source code, Lean~4 specifications, mutation scripts, raw results
(JSON), and the mutmut cache are included in the repository under
\texttt{studies/mutation\_testing/} and \texttt{formal/}.

\bibliographystyle{plain}
\bibliography{references}

\end{document}
